\documentclass[12pt, a4paper]{article}

% --- PREÁMBULO ---
\usepackage[utf8]{inputenc}
\usepackage[spanish]{babel}
\usepackage{geometry}
\usepackage{hyperref}
\usepackage{listings}
\usepackage{xcolor}
\usepackage{graphicx}
\usepackage{parskip} % Para tener espacio entre párrafos sin sangría
\usepackage{verbatim}

% Configuración de márgenes
\geometry{a4paper, total={170mm,257mm}, left=20mm, top=20mm}

% Configuración de hyperref para links
\hypersetup{
    colorlinks=true,
    linkcolor=blue,
    filecolor=magenta,      
    urlcolor=cyan,
    pdftitle={Manual y Bitácora de Conversaciones - Proyecto IA-BTC-FORECAST},
    pdfpagemode=FullScreen,
}

% --- Estilo para los bloques de código ---
\definecolor{codegreen}{rgb}{0,0.6,0}
\definecolor{codegray}{rgb}{0.5,0.5,0.5}
\definecolor{codepurple}{rgb}{0.58,0,0.82}
\definecolor{backcolour}{rgb}{0.95,0.95,0.92}

\lstdefinestyle{mystyle}{
    backgroundcolor=\color{backcolour},   
    commentstyle=\color{codegreen},
    keywordstyle=\color{magenta},
    numberstyle=\tiny\color{codegray},
    stringstyle=\color{codepurple},
    basicstyle=\ttfamily\footnotesize,
    breakatwhitespace=false,         
    breaklines=true,                 
    captionpos=b,                    
    keepspaces=true,                 
    numbers=left,
    literate={ó}{{\'o}}1 {á}{{\'a}}1 {é}{{\'e}}1 {í}{{\'i}}1 {ú}{{\'u}}1
             {Ó}{{\'O}}1 {Á}{{\'A}}1 {É}{{\'E}}1 {Í}{{\'I}}1 {Ú}{{\'U}}1
             {ñ}{{\~n}}1 {Ñ}{{\~N}}1 {¿}{{?`}}1 {¡}{{!`}}1,
    numbersep=5pt,                  
    showspaces=false,                
    showstringspaces=false,
    showtabs=false,                  
    tabsize=2
}
\lstset{style=mystyle}

% --- Comandos para diálogo ---
\newenvironment{usuario}{\par\noindent\textbf{Usuario:}\itshape\par\ignorespaces}{\par}
\newenvironment{gemini}{\par\noindent\textbf{Gemini:}\par\ignorespaces}{\par}

% --- DOCUMENTO ---
\begin{document}

\title{Manual y Bitácora de Conversaciones \\ Proyecto IA\_BTC\_FORECAST}
\author{CERO \& Gemini Code Assist}
\date{\today}
\maketitle

\begin{abstract}
    Este documento sirve como una bitácora de las conversaciones mantenidas con el asistente de IA Gemini durante el desarrollo del proyecto \texttt{IA\_BTC\_FORECAST}. Su propósito es documentar las preguntas, respuestas, decisiones de diseño y fragmentos de código generados.
\end{abstract}

\tableofcontents
\newpage

\section{Introducción a Google Colab}

\subsection{Conocimiento sobre la plataforma}

\begin{usuario}
Gemini, ¿conoces una plataforma llamada "Google Colab"?
\end{usuario}

\begin{gemini}
Sí, por supuesto. Google Colaboratory, o "Colab", es un servicio en la nube gratuito de Google que permite escribir y ejecutar código Python en el navegador. Es un entorno de Jupyter Notebook sin configuración, con acceso a hardware gratuito (GPUs y TPUs), ideal para Machine Learning.
\end{gemini}

\subsection{Soporte de Lenguajes}

\begin{usuario}
¿Esa plataforma solo sirve para programar en Python?
\end{usuario}

\begin{gemini}
Principalmente sí, está optimizado para Python. Sin embargo, también se pueden ejecutar comandos de shell (usando \texttt{`!`}) y, con algo de configuración, se pueden instalar kernels para otros lenguajes como R o Julia.
\end{gemini}

\subsection{Requisitos de Registro}

\begin{usuario}
¿Para usar esa plataforma hay que registrarse en la misma?
\end{usuario}

\begin{gemini}
Sí, pero no es un registro separado. Solo necesitas una cuenta de Google. Si ya usas Gmail o Google Drive, puedes usar esa misma cuenta. Tus notebooks se guardan directamente en tu Google Drive.
\end{gemini}

\section{Integración del Proyecto con Colab y LaTeX}

\subsection{Subir el proyecto a Colab usando Git}

\begin{usuario}
¿Cómo puedo subir mi proyecto \texttt{IA\_BTC\_FORECAST} a Google Colab usando Git?
\end{usuario}

\begin{gemini}
El proceso es sencillo:
\begin{enumerate}
    \item Clonar el repositorio: En una celda de Colab, ejecuta:
    
    \begin{lstlisting}[language=bash]
        !git clone https://github.com/ceroiglesias/IA_BTC_FORECAST.git
    \end{lstlisting}

    \item Cambiar de directorio: Usa el comando mágico:
    \begin{lstlisting}[language=bash]
%cd IA_BTC_FORECAST
    \end{lstlisting}
    \item Instalar dependencias: Ejecuta:
    \begin{lstlisting}[language=bash]
!pip install -r requirements.txt
    \end{lstlisting}
    \item Ejecutar scripts: Llama a tus scripts con:
    \begin{lstlisting}[language=bash]
!python nombre_script.py
    \end{lstlisting}
\end{enumerate}
\end{gemini}

\subsection{Creación de Guía en LaTeX}

\begin{usuario}
¿Podrías vaciarme toda esta información que me acabas 
de dar en un archivo latex?
\end{usuario}

\begin{gemini}
Claro. Generé un documento \texttt{.tex} completo con toda la información sobre Google Colab, sus características y la guía 
para subir tu proyecto. El archivo incluye un preámbulo bien configurado para formatear el texto y los bloques de código.
\end{gemini}

\subsection{Corrección de Error de Compilación en LaTeX}

\begin{usuario}
Tenemos un error al compilar el archivo latex. [Se adjunta log 
de error].
\end{usuario}

\begin{gemini}
El error \texttt{Invalid UTF-8 byte sequence} se debe a que el 
paquete \texttt{listings} no está interpretando correctamente 
los caracteres con tilde dentro de los comentarios del código. 
La solución es añadir la opción \texttt{literate} a la 
configuración de \texttt{lstdefinestyle} para que LaTeX 
reemplace los caracteres especiales (como 'ó', 'á', 'ñ') por 
sus comandos correspondientes 
(como \texttt{\{\'o\}}, \texttt{\{\'a\}}, \texttt{\{\~n\}}).
\end{gemini}

\section{Desarrollo y Herramientas}

\subsection{Revisión y Mejora del script \texttt{commit.sh}}

\begin{usuario}
Revisa el script "commit.sh" que te muestro en el contexto.
\end{usuario}

\begin{gemini}
El script es excelente. Propuse algunas mejoras para aumentar su robustez y claridad:
\begin{itemize}
    \item Usar un "here document" para la función de ayuda, haciéndola más mantenible.
    \item Hacer que el mensaje de commit sea obligatorio.
    \item Añadir una comprobación interactiva para \texttt{flake8}: si encuentra errores, pregunta al usuario si desea continuar.
    \item Comprobar si hay cambios pendientes antes de intentar hacer el commit para evitar errores.
\end{itemize}
Se proporcionó una versión refactorizada del script con estos cambios.
\end{gemini}

\subsection{Capacidades de Asistencia en IA con TensorFlow}

\begin{usuario}¿Estás en capacidad de ayudarme a desarrollar un 
script IA en python que utilize la librería Tensor Flow?
\end{usuario}

\begin{gemini}
¡Por supuesto que sí! Es una de mis especialidades. Puedo ayudarte a escribir código desde cero, explicar conceptos de TensorFlow/Keras, depurar tus scripts o trabajar contigo para optimizar tu modelo \texttt{IA\_BTC\_FORECAST}.
\end{gemini}

\section{Creación de esta Bitácora}

\usuario{Perfecto Gemini, pero para no perder esta y futuras converzaciones quiero que crees un archivo en formato latex con esta converzación y luego guardalo en el directorio "/home/cero/MEGA/VSCODE/Binance/Python/IA\_BTC\_FORECAST/Documentacion" con el nombre de "manual.tex. Mas adelante, cuando te pida que guardes una converzación, la anexas a ese archivo (manual.text) en formato latex.}

\begin{gemini}
¡Excelente idea! He creado este documento, \texttt{manual.tex}, en la ruta especificada. 
Contiene un resumen de nuestra conversación hasta la fecha. A partir de ahora, cuando me 
lo indiques, añadiré nuestras nuevas interacciones a este mismo archivo para mantener un 
registro completo.
\end{gemini}

\section{Análisis de la Arquitectura del Modelo}

\subsection{Explicación Detallada de Capas LSTM y Dense}

\begin{usuario}Ya compilé el archivo manual.tex sin errores. 
Ahora explícame en detalle las capas \texttt{LSTM} y 
\texttt{Dense} de mi modelo actual.
\end{usuario}

\begin{gemini}
    ¡Genial que el archivo \texttt{manual.tex} haya compilado 
    correctamente! Ahora tenemos una excelente bitácora para 
    nuestro trabajo.

Vamos a sumergirnos en el corazón de tu modelo de predicción. 
Las capas \texttt{LSTM} y \texttt{Dense} son los ladrillos 
fundamentales con los que estás construyendo tu red neuronal.

\subsubsection{La Capa \texttt{LSTM} (Long Short-Term Memory)}

Piensa en la capa \texttt{LSTM} como el \textbf{cerebro con memoria} de tu modelo. Su 
superpoder es entender \textbf{secuencias y contexto a lo largo del tiempo}. Esto es 
absolutamente crucial para tu proyecto, porque el precio futuro del Bitcoin no depende 
solo del precio de ayer, sino de la tendencia, el impulso y los patrones que se han 
formado durante los últimos días, semanas o meses.

\paragraph{¿Cómo funciona? (Analogía)}
Imagina que estás leyendo una novela. Para entender la trama, no solo lees la palabra actual, sino que recuerdas a los personajes y los eventos que ocurrieron en capítulos anteriores.
\begin{itemize}
    \item \textbf{RNN Simple (Red Neuronal Recurrente):} Sería como una persona con muy mala memoria a corto plazo. Recuerda la frase anterior, pero olvida rápidamente lo que pasó al principio del libro. Esto se conoce como el "problema del gradiente desvaneciente".
    \item \textbf{LSTM:} Es como un lector atento. Tiene un "cuaderno de notas" especial (el \textbf{estado de la celda} o \textit{cell state}) donde apunta las cosas importantes de la trama. Además, tiene tres "supervisores" o \textbf{compuertas} (\textit{gates}) que gestionan ese cuaderno:
    \begin{enumerate}
        \item \textbf{Compuerta de Olvido (Forget Gate):} Decide qué información del cuaderno ya no es relevante y debe ser borrada. Por ejemplo: "Este personaje secundario murió, ya no necesito recordarlo". En tu caso: "La pequeña fluctuación de precio de hace 30 días ya no es importante para la tendencia actual".
        \item \textbf{Compuerta de Entrada (Input Gate):} Decide qué nueva información de la palabra que está leyendo es lo suficientemente importante como para añadirla al cuaderno. Por ejemplo: "¡Un nuevo villano ha aparecido! Esto es clave, lo apunto". En tu caso: "El precio acaba de romper una resistencia importante, ¡esto es una señal clave y debo recordarla!".
        \item \textbf{Compuerta de Salida (Output Gate):} Decide qué parte del cuaderno es relevante para la acción inmediata que debe tomar (la predicción). Por ejemplo: "Basado en mis notas sobre el héroe y el villano, creo que el próximo capítulo será una pelea". En tu caso: "Basado en la tendencia alcista que he memorizado y la reciente ruptura de resistencia, mi predicción para el siguiente día será un valor más alto".
    \end{enumerate}
\end{itemize}

\paragraph{Parámetros Claves en tu Modelo}
\begin{itemize}
    \item \texttt{units=50}: Esto significa que tu capa LSTM tiene 50 "neuronas" o 
    "inidades de memoria". Cada una de estas unidades mantiene su propio 
    "cuaderno de notas" (su propio estado de celdas y estado oculto) para aprender
    un patrón o aspecto diferente de los datos. Más unidades pueden capturar patrones
    más complejos pero también aumentan el riesgo de sobreajuste (memorizar los
    datos de entrenamiento en lugar de aprender a generalizar).
    \item \texttt{return\_sequences=True}: Este es un parámetro vital cuando apilas 
    varias capas LSTM. Le dice a la capa: "No me des solo tu conclusión final después de 
    ver los 60 días; dame tus 'pensamientos' o estado oculto para cada uno de los 60 días". 
    Esto permite que la siguiente capa LSTM reciba una secuencia completa y pueda seguir 
    construyendo sobre los patrones que la capa anterior encontró. Si fuera la última capa 
    LSTM antes de una capa \texttt{Dense}, normalmente se pondría en \texttt{False}.
    \item \texttt{input\_shape=(X\_train.shape[1], 1)}: Esto es solo para la primera capa. 
    Le dice al modelo cómo es una muestra de entrada. En tu caso, es \texttt{(60, 1)}, lo que 
    significa que cada muestra es una secuencia de 60 pasos de tiempo (días) con 1 característica 
    por paso (el precio de cierre).
\end{itemize}

\subsubsection{La Capa \texttt{Dense}}
Después de que las capas \texttt{LSTM} han hecho su trabajo de entender la secuencia y
extraer patrones temporales, la capa \texttt{Dense} entra en acción. Piensa en ella como el 
\textbf{cerebro que toma decisiones}. A diferencia de la LSTM, no tiene memoria de secuencias. 
Su trabajo es tomar la información que le pasan (sin importar el orden) y encontrar las 
mejores combinaciones para llegar a una conclusión.

\paragraph{¿Cómo funciona?}
Cada neurona en una capa \texttt{Dense} está conectada a \textbf{todas} las neuronas de 
la capa anterior. Recibe todas las salidas, las pondera (les da más importancia a unas 
que a otras) y, basándose en esas ponderaciones, produce una salida.

\paragraph{Su Rol en tu Modelo}
En tu arquitectura, las capas \texttt{Dense} cumplen dos funciones:
\begin{enumerate}
    \item \textbf{Capas \texttt{Dense} Intermedias (\texttt{units=25}):} Estas capas actúan como un "procesador de conclusiones". La última capa LSTM le pasa un resumen muy rico y comprimido de todos los patrones temporales que encontró en los 60 días. La capa \texttt{Dense} toma este resumen y aprende a combinar las diferentes piezas de información. Es como si dijera: "Ok, la LSTM me dice que hay una tendencia alcista (patrón A) pero también una pequeña señal de reversión a corto plazo (patrón B). ¿Cómo combino A y B para refinar mi predicción?". Ayuda al modelo a aprender relaciones no lineales más complejas sobre los patrones ya extraídos.
    \item \textbf{Capa \texttt{Dense} de Salida (\texttt{units=1}):} Esta es la capa final, la que da la respuesta. Tiene una sola neurona porque tu objetivo es predecir un solo valor: el precio del Bitcoin para el siguiente día. Esta neurona toma toda la información procesada por las capas anteriores y la condensa en un único número, que es la predicción final del modelo.
\end{enumerate}

\subsubsection{Resumen del Flujo de Información}
\begin{enumerate}
    \item \textbf{Entrada:} Una ventana de 60 días de precios de BTC.
    \item \textbf{Capas LSTM:} Procesan la secuencia día a día, recordando patrones y 
    tendencias a largo plazo y pasando una secuencia enriquecida a la siguiente capa.
    \item \textbf{Capa \texttt{Dense(25)}:} Recibe el resumen final de las LSTM y lo 
    procesa para encontrar combinaciones complejas entre los patrones detectados.
    \item \textbf{Capa \texttt{Dense(1)}:} Recibe la información de la capa anterior y la 
    concentra en una única salida numérica: la predicción del precio para el día 61.
\end{enumerate}
Espero que esta explicación detallada te ayude a entender mejor cómo "piensa" tu modelo. ¡
Es una arquitectura muy potente para el problema que estás resolviendo!
\end{gemini}

\section{Flujo de Trabajo en Google Colab}

\subsection{Inicio de un Nuevo Notebook}

\begin{usuario}
Una vez que se entra a "Google Colab" hay que 
iniciar un nuevo notebook. Cierto?
\end{usuario}

\begin{gemini}
¡Correcto! Una vez que entras a Google Colab, lo más común es 
iniciar un nuevo notebook para empezar a trabajar.

Cuando visitas \href{https://colab.research.google.com}{colab.research.google.com}, generalmente te aparece una ventana emergente que te da varias opciones:
\begin{itemize}
    \item \textbf{Crear un notebook nuevo:} La opción más directa para un proyecto desde cero.
    \item \textbf{Abrir un notebook reciente:} Para continuar con tu trabajo.
    \item \textbf{Abrir desde Google Drive:} Si tienes notebooks guardados allí.
    \item \textbf{Abrir desde GitHub:} Puedes pegar la URL de un repositorio o un notebook específico.
    \item \textbf{Subir un archivo:} Si tienes un archivo \texttt{.ipynb} en tu ordenador.
\end{itemize}

Para nuestro proyecto \texttt{IA\_BTC\_FORECAST}, el flujo de trabajo ideal es exactamente ese:
\begin{enumerate}
    \item Iniciar un notebook nuevo.
    \item En la primera celda, clonar tu repositorio de GitHub con:
    \begin{lstlisting}[language=bash] !git clone ...
    \end{lstlisting}
\end{enumerate}

De esta manera, mantienes un entorno limpio y reproducible cada vez que trabajas.
\end{gemini}


\subsection{Clonado de Repositorios (Público vs. Privado)}

\begin{usuario}
Creo que para clonar un repositorio de github, antes hay que 
hacerlo público?
\end{usuario}

\begin{gemini}
¡Esa es una excelente pregunta! No, no es estrictamente 
necesario. Puedes clonar tanto repositorios públicos como 
privados, pero el procedimiento es diferente:

\begin{enumerate}
    \item \textbf{Repositorios Públicos:}
    \begin{itemize}
        \item Cualquiera puede clonarlos sin necesidad de autenticarse.
        \item Este es el caso de tu proyecto \texttt{IA\_BTC\_FORECAST}, por eso el comando:
        \begin{lstlisting}[language=bash] !git clone https://...
        \end{lstlisting}
        funciona directamente en Colab.
    \end{itemize}

    \item \textbf{Repositorios Privados:}
    \begin{itemize}
        \item Necesitas tener permisos de acceso y debes autenticarte.
        \item Las formas más comunes son usando un \textbf{Token de Acceso Personal (PAT)} con la URL HTTPS, o configurando una 
        \textbf{clave SSH}.
        \item Por ejemplo, para clonar usando un PAT, la URL se vería así:
        \begin{lstlisting}[language=bash]
!git clone https://<TOKEN>@github.com/usuario/repositorio.git
        \end{lstlisting}
        \item Alternativamente, puedes configurar una clave SSH en Colab, pero esto es un poco más avanzado y requiere pasos 
        adicionales.
    \end{itemize}
\end{enumerate}

En resumen, para tu proyecto actual no tienes que preocuparte porque es público. Pero si en el futuro trabajas con un repositorio privado, solo necesitarás autenticarte para poder clonarlo.
\end{gemini}

\subsection{Error Común: URL de Navegación vs. URL de Clonado}

\begin{usuario}
Le doy el siguiente camino: 
https://github.com/cero97-ctrl/repositorio/tree/main/Binance/Python. 
Pero me dice que no encuentra el repositorio.
\end{usuario}


\begin{gemini}
¡Entiendo perfectamente el problema! Es un error muy común. La 
URL que estás usando es la URL de \textbf{navegar} por una 
subcarpeta en la web de GitHub, no para clonar.

El comando \texttt{git clone} necesita la URL principal del 
repositorio, que siempre descarga el proyecto completo.

\paragraph{La Solución}
\begin{enumerate}
    \item \textbf{Usa la URL correcta para clonar.} Para tu repositorio, es:
    \begin{lstlisting}[language=bash]
https://github.com/cero97-ctrl/repositorio.git
    \end{lstlisting}
    \item \textbf{Clona el repositorio completo:}
    \begin{lstlisting}[language=bash]
!git clone https://github.com/cero97-ctrl/repositorio.git
    \end{lstlisting}
    \item \textbf{Navega a la carpeta deseada:} Una vez clonado, puedes moverte a la 
    subcarpeta que necesitas:
\end{enumerate}
\begin{lstlisting}[language=Python]
import os
os.chdir('repositorio/Binance/Python')
\end{lstlisting}
\end{gemini}

\subsection{Visualización del Contenido del Repositorio}

\begin{usuario}
¿Cómo muestro el contenido del repositorio?
\end{usuario}

\begin{gemini}
    ¡Claro! Una vez que has clonado el repositorio, este se descarga en una nueva carpeta. 
    Para ver su contenido desde Colab, puedes usar el comando \texttt{ls} de Linux, 
    prefijado con un signo de exclamación (\texttt{!}).

    \paragraph{Pasos para ver el contenido}
    \begin{enumerate}
        \item \textbf{Clona el repositorio (si no lo has hecho ya):}
        \begin{lstlisting}[language=bash]
            !git clone https://github.com/cero97-ctrl/repositorio.git
        \end{lstlisting}
        \item \textbf{Lista el contenido de la nueva carpeta:}
        \begin{lstlisting}[language=bash]
            !ls repositorio
        \end{lstlisting}
        \item \textbf{Para una vista más detallada} (con permisos, tamaños y archivos ocultos):
        \begin{lstlisting}[language=bash]
!ls -lha repositorio
        \end{lstlisting}
    \end{enumerate}

    \paragraph{Flujo de trabajo alternativo}
        También puedes 'entrar' en el directorio del repositorio con:
        \begin{lstlisting}[language=bash] 
            %cd repositorio
        \end{lstlisting}
        y luego ejecutar \texttt{!ls -lha} para listar el contenido del directorio actual. Esto es útil si vas a ejecutar varios 
        comandos dentro de esa carpeta.
\end{gemini}


\subsection{Activación y Verificación de GPU}

\begin{usuario}
¿Cómo puedo activar y verificar el uso de la GPU en este 
notebook de Colab?
\end{usuario}

\begin{gemini}
¡Claro! Activar la GPU en Colab es uno de sus puntos más 
fuertes y es muy sencillo. 
Aquí te explico cómo hacerlo y cómo comprobar que está 
funcionando

\subsubsection{Paso 1: Activar la GPU}
El entorno de ejecución de Colab usa una CPU por defecto. Para cambiar a una GPU, sigue 
estos pasos:
\begin{enumerate}
    \item Ve al menú superior y haz clic en \textbf{Entorno de ejecución} (Runtime).
    \item En el menú desplegable, selecciona \textbf{Cambiar tipo de entorno de ejecución} 
    (Change runtime type).
    \item Se abrirá una ventana. Busca el menú desplegable llamado 
    \textbf{Acelerador por hardware} (Hardware accelerator).
    \item Selecciona \textbf{T4 GPU} (o la opción de GPU que aparezca).
    \item Haz clic en \textbf{Guardar} (Save).
\end{enumerate}
\textbf{Importante:} Este cambio reiniciará tu entorno de ejecución, por lo que tendrás que 
volver a ejecutar las celdas anteriores (como la de clonar el repositorio o instalar 
dependencias).

\subsubsection{Paso 2: Verificar que la GPU está activa}
Una vez que el entorno se haya reiniciado, puedes verificar que la GPU está disponible de 
dos maneras muy sencillas.

\paragraph{Método 1: Con un comando de shell}
Este es el método más rápido y común. Ejecuta la siguiente celda:
\begin{lstlisting}[language=bash]
!nvidia-smi
\end{lstlisting}
Si la GPU está activa, verás una tabla con información detallada sobre la GPU asignada 
(generalmente una Tesla T4). Si obtienes un error como \texttt{nvidia-smi: command not found}, 
significa que la GPU no está activada.

\paragraph{Método 2: Con TensorFlow}
También puedes pedirle a TensorFlow que liste los dispositivos GPU que detecta. Ejecuta
esta celda:
\begin{lstlisting}[language=Python]
import tensorflow as tf
print(tf.config.list_physical_devices('GPU'))
\end{lstlisting}
Si la GPU está funcionando, la salida será una lista que contiene información sobre el 
dispositivo GPU, algo como: \texttt{[PhysicalDevice(name='/physical\_device:GPU:0', 
device\_type='GPU')]}. Si la lista está vacía (\texttt{[]}), TensorFlow no está detectando 
la GPU.
\end{gemini}

\end{document}